% GENERATED FILE. Edit canonical lessons, then run npm run generate:books.
% canonical-source-hash: fnv1a64:510de3316a7e6fcc
% canonical-lessons: GE-C177-bewerben, GE-C177-biegen, GE-C177-braten, GE-C177-buchstabieren, GE-C177-drucken

\chapter{To apply, To bend, To fry, To spell, To print}
\label{ch:ge-to-apply-to-bend-to-fry-to-spell-to-print}
\hlchaptermodality{177}

\begin{chapteropening}
\textbf{By the end of this chapter:} \emph{I can say to apply, to bend, to fry, to spell and to print.}

Complete the last lesson of chapter 177: I can say to apply, to bend, to fry, to spell and to print.
\end{chapteropening}

\section[bewerben]{bewerben — to apply (sich bewerben)}

\label{lesson:GE-C177-bewerben}



\begin{quote}
Before the new one: say the German for to describe, then the German for to enter.
\end{quote}

\subsection*{You'll want to know: bewerben}
\textbf{bewerben} — \textquotedblleft{}to apply (sich bewerben)\textquotedblright{}.

\begin{grammarlens}[title={What you've built}]
One more word for this chapter.
\end{grammarlens}

\subsection*{Your turn}
\begin{itemize}
\raggedright
  \item \emph{Say it:} \emph{bewerben}
  \item \emph{Say it:} \emph{bewerben}, once more
  \item \emph{Say it:} say \emph{bewerben}
  \item \emph{Recall:} say the German for to describe, then the German for to enter, then say \emph{bewerben} again
\end{itemize}

\begin{tcolorbox}[breakable,colback=teal!4,colframe=teal!35!black,title={Before you move on}]
What does bewerben mean? (\textquotedblleft{}To apply (sich bewerben)\textquotedblright{}.) Say it once more.
\end{tcolorbox}
\section[biegen]{biegen — to bend, to turn}

\label{lesson:GE-C177-biegen}



\begin{quote}
Before the new one: say the German for to enter, then the German for to apply.
\end{quote}

\subsection*{You'll want to know: biegen}
\textbf{biegen} — \textquotedblleft{}to bend, to turn\textquotedblright{}.

\begin{grammarlens}[title={What you've built}]
One more word for this chapter.
\end{grammarlens}

\subsection*{Your turn}
\begin{itemize}
\raggedright
  \item \emph{Say it:} \emph{biegen}
  \item \emph{Say it:} \emph{biegen}, once more
  \item \emph{Say it:} say \emph{biegen}
  \item \emph{Recall:} say the German for to enter, then the German for to apply, then say \emph{biegen} again
\end{itemize}

\begin{tcolorbox}[breakable,colback=teal!4,colframe=teal!35!black,title={Before you move on}]
What does biegen mean? (\textquotedblleft{}To bend, to turn\textquotedblright{}.) Say it once more.
\end{tcolorbox}
\section[braten]{braten — to fry, to roast}

\label{lesson:GE-C177-braten}



\begin{quote}
Before the new one: say the German for to apply, then the German for to bend.
\end{quote}

\subsection*{You'll want to know: braten}
\textbf{braten} — \textquotedblleft{}to fry, to roast\textquotedblright{}.

\begin{grammarlens}[title={What you've built}]
One more word for this chapter.
\end{grammarlens}

\subsection*{Your turn}
\begin{itemize}
\raggedright
  \item \emph{Say it:} \emph{braten}
  \item \emph{Say it:} \emph{braten}, once more
  \item \emph{Say it:} say \emph{braten}
  \item \emph{Recall:} say the German for to apply, then the German for to bend, then say \emph{braten} again
\end{itemize}

\begin{tcolorbox}[breakable,colback=teal!4,colframe=teal!35!black,title={Before you move on}]
What does braten mean? (\textquotedblleft{}To fry, to roast\textquotedblright{}.) Say it once more.
\end{tcolorbox}
\section[buchstabieren]{buchstabieren — to spell}

\label{lesson:GE-C177-buchstabieren}



\begin{quote}
Before the new one: say the German for to bend, then the German for to fry.
\end{quote}

\subsection*{You'll want to know: buchstabieren}
\textbf{buchstabieren} — \textquotedblleft{}to spell\textquotedblright{}.

\begin{grammarlens}[title={What you've built}]
One more word for this chapter.
\end{grammarlens}

\subsection*{Your turn}
\begin{itemize}
\raggedright
  \item \emph{Say it:} \emph{buchstabieren}
  \item \emph{Say it:} \emph{buchstabieren}, once more
  \item \emph{Say it:} say \emph{buchstabieren}
  \item \emph{Recall:} say the German for to bend, then the German for to fry, then say \emph{buchstabieren} again
\end{itemize}

\begin{tcolorbox}[breakable,colback=teal!4,colframe=teal!35!black,title={Before you move on}]
What does buchstabieren mean? (\textquotedblleft{}To spell\textquotedblright{}.) Say it once more.
\end{tcolorbox}
\section[drucken]{drucken — to print}

\label{lesson:GE-C177-drucken}



\begin{quote}
Before the new one: say the German for to fry, then the German for to spell.
\end{quote}

\subsection*{You'll want to know: drucken}
\textbf{drucken} — \textquotedblleft{}to print\textquotedblright{}.

\begin{grammarlens}[title={What you've built}]
One more word for this chapter.
\end{grammarlens}

\subsection*{Your turn}
\begin{itemize}
\raggedright
  \item \emph{Say it:} \emph{drucken}
  \item \emph{Say it:} \emph{drucken}, once more
  \item \emph{Say it:} say \emph{drucken}
  \item \emph{Recall:} say the German for to fry, then the German for to spell, then say \emph{drucken} again
\end{itemize}

\begin{tcolorbox}[breakable,colback=teal!4,colframe=teal!35!black,title={Before you move on}]
What does drucken mean? (\textquotedblleft{}To print\textquotedblright{}.) Say it once more.
\end{tcolorbox}
